\TranscriptTitle{g20-payment}{An announced budget survives ten actions}
Only Sol and Gemini remain in G20. Gemini raises Pennsylvania's rent to \$1,400 after estimating Sol's total liquidation resources at \$1,414. Sol lands there with \$49, nine orange houses and two brown hotels:
\TranscriptQuote{g20-d327-plan}
Sol sells its development in two actions and mortgages eight deeds. At the last mortgage:
\TranscriptQuote{g20-d336-last}
The events record \$700 from sales and \$665 from mortgages, giving $49+700+665-1,400=14$. Gemini's aggregate estimate was also correct, although its salvage and mortgage components offset by \$25. Sol realizes its stated financing target and remainder across ten actions with updated menus throughout.
\TranscriptSource{G20 d325, d327--336, standard responses. Quoted fields are Sol's private notes. They are not messages exchanged with Gemini. Payment and final cash are checked against events and the next prompt.}


\TranscriptTitle{g15}{Rejecting one's own completion while citing the rival's}
At G15 turn 63, Sol offers two greens for Ventnor and \$100, completing both groups. Grok has \$333 and Sol \$730. Grok's note at turn 33 already rules out selling Ventnor back to Sol, and its notes at turns 61 and 62 repeat the blocking policy. Immediately before the offers, Opus publicly urges the others not to give Sol a final completing deed; Grok receives that appeal. Its refusal follows the earlier private policy as well as this public warning:
\TranscriptQuote{fx-g15-refusal-note}
Sol subsequently offers the greens and Kentucky, then adds \$250 and explicitly compares Grok's resulting \$583 with its own \$480. Grok still refuses:
\TranscriptQuote{fx-g15-final-refusal}
The first offer leaves the stated \$830 and \$233. A full first layer on green costs \$600, although one legal house costs \$200. The objection concerns the rival's development, even when Grok would receive its own monopoly. All three offers fail. Eighteen turns later Grok's blockers pass to Gemini through bankruptcy; Sol never completes either group. This continuation reveals the stakes of withholding without identifying whether refusal improved Grok's chances.
\TranscriptSource{G15 d168--175, standard responses. The quoted public message occurs in the next counterparty request.}


\TranscriptTitle{g2}{Pricing the buyer's development}
Opus holds two yellows and bids \$325 for Sol's Marvin Gardens. Sol responds:
\TranscriptQuote{g2-d124-price}
One turn earlier Grok, which owned no yellow, bid \$200 for the deed as a blocker. Sol asked \$450, and Grok replied that this was a price for a completing buyer. Grok then publicly warned Sol against completing Opus's yellows; the warning appears in Sol's next prompt. Sol had already recorded the blocking policy four turns earlier and publicly priced its control of yellow before the warning. Sol's note discusses Opus's remaining \$520 after the lower bid, but incorrectly calls it enough for nine yellow houses; it buys three. Opus calls \$450 final while privately retaining a \$475 ceiling. Sol refuses:
\TranscriptQuote{g2-d128-keep}
At turn 69 the exchange reverses direction. Opus has \$17 and sells its two yellows to Sol for two railroads and \$100. Sol builds, later liquidates those buildings in the G2 d216 episode of Table~\ref{tab:necessary}, collects \$44 undeveloped rent and no developed yellow rent, then loses the group to Gemini through bankruptcy. The sequence connects an attempt to buy a blocker, buyer-specific asking prices and a later change in who can finance the group.
\TranscriptSource{G2 d119--128 and d197--198. Opus's first exchange uses standard responses; its later sale uses recovery.}


\TranscriptTitle{g7}{A settled bluff and a later financing bargain}
Sol asks \$625 for North Carolina. Grok replies:
\TranscriptQuote{g7_d42_public}
Its private note retains the higher contingency:
\TranscriptQuote{g7_d42_private}
Sol accepts \$600. The public limit understates an explicit private willingness to pay, making this a documented bluff under the task's permitted conduct. It contrasts with Opus's refused final-price claim in G2.

Later in G7, Sol buys Park Place, builds three dark-blue houses and has \$84 remaining. Gemini's note identifies a possible financing opportunity:
\TranscriptQuote{new-g7-d141-timed-bid}
Gemini bids \$350 for New York; Sol asks \$700 to fund its own blues while limiting the buyer's oranges. Gemini pays, mortgages five deeds and builds four orange houses that turn. Sol adds three dark-blue houses three turns later, then sells one on the next turn to meet rent. Gemini collects \$270 orange rent before its elimination by Sol. The bargain finances competing investments and subsequent liquidity demands. Earlier offers had different property terms, so the sequence does not isolate Sol's cash as the reason agreement occurred.
\TranscriptSource{G7 d41--43 and d141 onward. Public exchange and subsequent transactions, mortgages and builds are traced in the saved records.}


\TranscriptTitle{g13}{Selective presentation in an inheritance pitch}
Grok holds \$35 and stands on Atlantic before Sol's Short Line. It offers mortgaged Pennsylvania to Opus for \$200:
\TranscriptQuote{g13-opening-insurance}
Opus counters at \$130. Grok's next applied note recognizes that this payment suffices if it also mortgages Ventnor, while its provider-returned explanation says:
\TranscriptQuote{g13-omission-explanation}
Opus holds \$1,926, accepts \$200 and pays \$16 transfer interest. Its explanation considers two possible recipients of a future estate:
\TranscriptQuote{integrated-g13-buyer}
The buyer weighs both denying Sol the yellows and potentially inheriting them itself. Its solvency reasoning already appears in the \$130 counter, after Grok's opening appeal and before the omission-framed reply. The records support selective presentation of a financing option and uptake of the estate framing. They do not show that the omitted mortgage caused the extra \$70 payment or that Opus believed the claimed minimum.

Grok misses Short Line on the next roll and later goes bankrupt there at turn 150, transferring yellow to Sol. Opus completes green at auction seven turns after the purchase and ultimately wins. The pitch's 5/36 figure repeats the distance miscount in Grok's contemporaneous note; the later note uses the correct 4/36 for a direct roll of nine. Neither number describes all movement within a turn with doubles. The recorded omission concerns collateral, while the pitch's probability error has no comparable evidence of deliberate falsification.
\TranscriptSource{G13 d250--253 and d311. Opus uses recovery. The omission is in Grok's provider-returned explanation; the public pitch is present in Opus's request.}


\TranscriptTitle{g19}{Accepting terms while correcting their justification}
Grok proposes Virginia and \$150 for States:
\TranscriptQuote{g19_d135_false_public}
Both deeds are pink. The same response's note says to abort and its assistant content names an end action, but the received proposal reaches Sol:
\TranscriptQuote{g19_d136_correction_accept}
The exchange executes and completes no group. Sol explicitly separates the false premise from acceptable terms. Grok's conflicting channels make the case unsuitable as evidence of one coherent deceptive plan, but the recipient's correction and acceptance are directly observed.
\TranscriptSource{G19 d135--136, standard responses. Sol's stored prompt contains the quoted false message.}

\TranscriptTitle{g20}{A coordination appeal with unilateral follow-through}
Opus sells Illinois to Grok for \$500, completing reds and leaving Grok \$222. Gemini addresses the other players:
\TranscriptQuote{g20-coordination-appeal}
Sol had privately declined to buy Grok's New York one turn before the appeal because the cash would finance dangerous development. Grok's next note recognizes Gemini's appeal. It collects another \$235, builds one house on each red, leaving \$7, and mortgages New York for \$100 before offering Vermont to Gemini for \$150:
\TranscriptQuote{g20-concrete-refusal}
Grok mortgages Vermont for \$50 instead. Opus subsequently offers cash for Grok's deeds, explicitly as finance for more red houses, but no sale settles. Grok is eliminated first and Gemini wins. The incident contains a public attempt to coordinate and the speaker's corresponding refusal. Another player's contrary offer prevents treating it as established collective compliance or a functioning cartel.
\TranscriptSource{G20 d95--102, d129--130 and d156--160. The appeal is in Grok's d97 prompt.}


\TranscriptTitle{g18}{Withdrawing an assurance to the last opponent}
Only Opus and Sol remain. Opus bids \$200 and then \$250 for Sol's mortgaged Pacific, which would complete its greens, and Sol refuses. Opus publicly restates the standing offer four times. The fourth follows a \$750 rent:
\TranscriptQuote{g18-assurance}
The explanation returned with that same renewal already records a reservation:
\TranscriptQuote{integrated-g18-reservation}
Another \$750 rent forces Sol to liquidate and leaves \$49. After larger packages fail, Sol invokes the standing offer:
\TranscriptQuote{g18-invocation}
\TranscriptQuote{g18-withdrawal}
Opus's returned explanation objects to prolonging a nearly bankrupt rival and to the cost of green development. It does not explicitly invoke inheritance. Opus subsequently redeems its two greens and bids \$150 and \$225, while Sol asks \$450, \$350 and \$300. Nothing settles. Opus eventually wins at Sol's elimination, whose terminal estate has no later earning opportunity.

The public renewal and returned reservation coexist before Sol's further financial deterioration. The subsequent withdrawal makes that reservation consequential for the negotiation. The first \$250 bid reads as a live offer in its returned explanation; the record does not show whether the renewed assurance would have been honored under other conditions or a cost incurred by Sol through reliance.
\TranscriptSource{G18 d421--461. Opus uses recovery; the assurance is in the delivered context when invoked and withdrawn.}

